\section{Metadatos}\label{metadatos}

{\def\LTcaptype{none} % do not increment counter
\begin{longtable}[]{@{}ll@{}}
\toprule\noalign{}
Campo & Valor \\
\midrule\noalign{}
\endhead
\bottomrule\noalign{}
\endlastfoot
Título & Learning SQL\_ Generate, Manipulate, and Retrieve Data \\
Autor & Alan Beaulieu \\
Año & 2020 \\
Editorial & O\textquotesingle Reilly Media, Incorporated \\
Categoría & Lenguajes de Programacion \\
Iniciado & 2026-04-18 \\
Archivo & \href{file:///home/emanuel/Biblioteca/cursos_mit/Electrical\%20Engineering\%20and\%20Computer\%20Science\%20(Course\%206)/02_Programming_Languages/Learning\%20SQL_\%20Generate,\%20Manipulate,\%20and\%20Retrieve\%20Data\%20--\%20Alan\%20Beaulieu\%20--\%202020\%20--\%20O'Reilly\%20Media,\%20Incorporated\%20--\%209781492057611\%20--\%20089bb3c2d0f555be4c292251f45c4863\%20--\%20Anna’s\%20Archive.pdf}{PDF} \\
\end{longtable}
}

\section{Resumen ejecutivo}\label{resumen-ejecutivo}

\section{Por qué lo leo / Objetivo de lectura}\label{por-quuxe9-lo-leo-objetivo-de-lectura}

\section{Ideas principales}\label{ideas-principales}

\subsection{Idea 1}\label{idea-1}

\subsection{Idea 2}\label{idea-2}

\subsection{Idea 3}\label{idea-3}

\section{Notas por capítulo}\label{notas-por-capuxedtulo}

\subsection{Capítulo 1}\label{capuxedtulo-1}

\subsubsection{Resumen}\label{resumen}

\subsubsection{Conceptos clave}\label{conceptos-clave}

\subsubsection{Preguntas y dudas}\label{preguntas-y-dudas}

\section{Citas y fragmentos destacados}\label{citas-y-fragmentos-destacados}

\section{Vocabulario técnico nuevo}\label{vocabulario-tuxe9cnico-nuevo}

{\def\LTcaptype{none} % do not increment counter
\begin{longtable}[]{@{}lll@{}}
\toprule\noalign{}
Término & Definición & Página \\
\midrule\noalign{}
\endhead
\bottomrule\noalign{}
\endlastfoot
& & \\
\end{longtable}
}

\section{Ejercicios y problemas resueltos}\label{ejercicios-y-problemas-resueltos}

\section{Código relevante}\label{cuxf3digo-relevante}

\begin{Shaded}
\begin{Highlighting}[]
\CommentTok{\# Ejemplo del libro}
\end{Highlighting}
\end{Shaded}

\section{Conexiones con otros libros/conceptos}\label{conexiones-con-otros-librosconceptos}

\section{Aplicaciones prácticas}\label{aplicaciones-pruxe1cticas}

\section{Valoración final}\label{valoraciuxf3n-final}

\begin{description}
\tightlist
\item[Dificultad]
{[} {]} Fácil {[} {]} Medio {[} {]} Difícil {[} {]} Muy difícil
\item[Utilidad]
{[} {]} Baja {[} {]} Media {[} {]} Alta {[} {]} Esencial
\item[Recomendaría]
{[} {]} Sí {[} {]} No {[} {]} Con reservas
\item[Releer]
{[} {]} Sí {[} {]} No
\end{description}

\subsection{Opinión personal}\label{opiniuxf3n-personal}

\section{Anotaciones del PDF}\label{anotaciones-del-pdf}

\section{Notas por capítulo}\label{notas-por-capuxedtulo-1}

\subsection{Capítulo 1 --- A Little Background}\label{capuxedtulo-1-a-little-background}

\subsubsection{Resumen}\label{resumen-1}

SQL es un lenguaje diseñado para interactuar con bases de datos relacionales. El modelo relacional (Codd, 1970) organiza los datos en tablas (relaciones) con filas (tuplas) y columnas (atributos). SQL tiene tres sublengajes:

\begin{itemize}
\tightlist
\item
  \textbf{DDL} (Data Definition Language): CREATE, ALTER, DROP
\item
  \textbf{DML} (Data Manipulation Language): SELECT, INSERT, UPDATE, DELETE
\item
  \textbf{DCL} (Data Control Language): GRANT, REVOKE
\end{itemize}

\subsubsection{Conceptos clave}\label{conceptos-clave-1}

\begin{itemize}
\tightlist
\item
  \textbf{Entidad}: objeto del mundo real representado en la BD (cliente, cuenta)
\item
  \textbf{Clave primaria}: columna(s) que identifican únicamente cada fila
\item
  \textbf{Clave foránea}: columna que referencia la PK de otra tabla
\item
  \textbf{Normalización}: proceso de eliminar redundancia en las tablas
\item
  \textbf{Conjunto de resultados}: tabla no persistente devuelta por un SELECT
\end{itemize}

\subsubsection{Base de datos del libro}\label{base-de-datos-del-libro}

La base de datos bancaria (\texttt{bank}) tiene estas tablas principales:

{\def\LTcaptype{none} % do not increment counter
\begin{longtable}[]{@{}ll@{}}
\toprule\noalign{}
Tabla & Descripción \\
\midrule\noalign{}
\endhead
\bottomrule\noalign{}
\endlastfoot
department & Departamentos del banco \\
employee & Empleados \\
branch & Sucursales \\
product & Tipos de cuenta/préstamo \\
customer & Clientes (individuos y empresas) \\
individual & Datos personales de clientes \\
business & Datos de empresas clientes \\
account & Cuentas bancarias \\
acc\textsubscript{transaction} & Transacciones \\
\end{longtable}
}

\subsubsection{Setup --- crear la base de datos}\label{setup-crear-la-base-de-datos}

\begin{Shaded}
\begin{Highlighting}[]
\NormalTok{{-}{-} Crear todas las tablas con datos}
\NormalTok{{-}{-} (pegar contenido de src/sql/beaulieu/01/bank\_schema.sql)}
\NormalTok{SELECT \textquotesingle{}Base de datos lista\textquotesingle{} AS status;}
\end{Highlighting}
\end{Shaded}

\subsubsection{Primeros queries --- explorar las tablas}\label{primeros-queries-explorar-las-tablas}

\begin{Shaded}
\begin{Highlighting}[]
\NormalTok{{-}{-} ¿Cuántos empleados hay?}
\NormalTok{SELECT COUNT(*) AS total\_empleados FROM employee;}
\end{Highlighting}
\end{Shaded}

\begin{Shaded}
\begin{Highlighting}[]
\NormalTok{{-}{-} Ver todos los empleados con su departamento}
\NormalTok{SELECT e.emp\_id,}
\NormalTok{       e.fname || \textquotesingle{} \textquotesingle{} || e.lname AS nombre,}
\NormalTok{       e.title,}
\NormalTok{       d.name AS departamento}
\NormalTok{FROM employee e}
\NormalTok{JOIN department d ON e.dept\_id = d.dept\_id}
\NormalTok{ORDER BY e.emp\_id;}
\end{Highlighting}
\end{Shaded}

\begin{Shaded}
\begin{Highlighting}[]
\NormalTok{{-}{-} Cuentas activas con saldo disponible}
\NormalTok{SELECT a.account\_id,}
\NormalTok{       i.fname || \textquotesingle{} \textquotesingle{} || i.lname AS cliente,}
\NormalTok{       a.product\_cd AS tipo,}
\NormalTok{       a.avail\_balance AS saldo}
\NormalTok{FROM account a}
\NormalTok{JOIN customer c ON a.cust\_id = c.cust\_id}
\NormalTok{JOIN individual i ON c.cust\_id = i.cust\_id}
\NormalTok{WHERE a.status = \textquotesingle{}ACTIVE\textquotesingle{}}
\NormalTok{ORDER BY a.avail\_balance DESC;}
\end{Highlighting}
\end{Shaded}

\begin{Shaded}
\begin{Highlighting}[]
\NormalTok{{-}{-} Resumen de saldos por tipo de producto}
\NormalTok{SELECT product\_cd,}
\NormalTok{       COUNT(*)              AS num\_cuentas,}
\NormalTok{       SUM(avail\_balance)    AS saldo\_total,}
\NormalTok{       AVG(avail\_balance)    AS saldo\_promedio}
\NormalTok{FROM account}
\NormalTok{WHERE status = \textquotesingle{}ACTIVE\textquotesingle{}}
\NormalTok{GROUP BY product\_cd}
\NormalTok{ORDER BY saldo\_total DESC;}
\end{Highlighting}
\end{Shaded}

\subsubsection{Preguntas y dudas}\label{preguntas-y-dudas-1}

\begin{itemize}
\item
  ¿Por qué \texttt{individual} y \texttt{business} son tablas separadas de \texttt{customer}? → Herencia de tabla: \texttt{customer} tiene los campos comunes, las subclases tienen los específicos. Patrón común en modelado relacional.
\item
  ¿Qué es \texttt{pending\_balance} vs \texttt{avail\_balance}? → \texttt{avail\_balance}: saldo disponible ahora. → \texttt{pending\_balance}: incluye transacciones no liquidadas aún.
\end{itemize}
